\scanpage{624}
\mspage{\uline{12}}
On a alors, pour $u:I\longrightarrow I'$

\[
\left\{\begin{aligned}
&u^*(\Phi(I'))\subset\deleted{\unsure{\Phi}}\Phi(I)\\
&u_*(\Psi[I])\subset\Psi[I']
\end{aligned}\right.\tag*{(1.22)}
\]

Il se pose encore la question si le couple
$(\Phi(I),\Psi[I])$ satisfait au \deleted{le} th. de factorisation,
sous des hypothèses convenables sur $(\Phi,\Psi)$.
Je peux seulement hasarder la

\uline{Question-conjecture} \uline{1.11}. Supposons que
\deleted{$(\Phi,\unsure{\Psi}$} $(\Phi,\Psi)$ soit un couple de Quillen
\uline{spécial}, en est-il de $=$ de $(\Phi(I),\Psi[I])$
pour \unsure{tt} $I$ dans $(\mathrm{Cat})$ ? \deleted{\unsure{Je co}} En d'autres
termes, si \deleted{$\Phi$ \unsure{est une}} $M$ est spécial
(condition (1.10 a)), et si $\Phi$ est une
partie $Q$-close à g. \uline{spéciale} de $M$ i.e.
admettant une \deleted{petite} partie $Q$-génératrice,
est-il vrai que pour \unsure{tte} $I$ dans $\mathrm{Cat}$,
$\Phi(I)$ est également spéciale ?
\inserted{Je conjecture que oui, en faisant au \unsure{besoin} des hypothèses}
% slanted
\margin{\unsure{suppl.} sur $M$ (existence d'une filtration \unsure{``cardinale''} ?)}
P. ex. si $\Phi_0$
est $Q$-génératrice dans $\Phi$, et \deleted{petite}, alors
$\Phi_0(I)$ est-elle $Q$-génératrice \unsure{dans} $\Phi(I)$ ?

\uline{Prop. 1.12} Soit $\Phi\longleftrightarrow\Psi$ dans $M$, soit $u:I\longrightarrow I'$
une flèche de $\mathrm{Cat}$, d'où

\diagram{\begin{tikzcd}[ampersand replacement=\&,column sep=large]
\text{(1.12.1)} \& M(I)\arrow[r,shift left=2.5,"u_!"]\arrow[r,shift right=2.5,"u_*"'] \& M(I')\arrow[l,"u^*"] .
\end{tikzcd}}

% slanted
\margin{cf \uline{$\overline{\mathrm{XVII}}$} 6.4 (p. \unsure{117|217|17}), mais où les notations pour $\Phi\langle I\rangle,\Psi(I)$ sont inversées.}
On a alors (1.19) \unsure{sc.|sic.|i.e.}

\[
\left\{\begin{aligned}
&u_!\Phi(\langle I\rangle)\subset\Phi\langle I'\rangle\\
&u^*(\Psi(\unsure{I|I'}))\subset\Psi(I'),
\end{aligned}\right.\tag*{(1.13.1)}
\]

et de plus l'équivalence

\[
u^*(\Phi\langle I'\rangle)\subset\Phi\langle I\rangle\Longleftrightarrow
u_*(\Psi(I))\subset\Psi(I').\tag*{(1.13.2)}
\]

Ces relations équiv. sont satisfaites dans chacun des 3 cas :
$1^\circ)$ $u$ est coétale
$2^\circ)$ $I$ est discrète
$3^\circ)$ $\forall i'$ dans $I'$, \deleted{toute} \inserted{\unsure{chaq.}} \unsure{comp. conn.} de $I_{/i'}$ admet un \unsure{objet} final.
% slanted
\margin{\uline{NB} On a $1^\circ\Rightarrow3^\circ$, $2^\circ\Rightarrow3^\circ$}
% archive sheet number: 623
